\RequirePackage{pdfmanagement-testphase} % tranparent package doesn't work without this
\DeclareDocumentMetadata{} % and this
\documentclass[a4paper, 11pt, oneside]{impression}

\hwnumber{1} % Homework number in the cover page
\title{Lista de Exercícios} % Title
\subtitle{Relatividade Geral - A Teoria da Gravitação de Einstein} % Subtitle
\author{Lucas R. Ximenes} % Author
\extra{11917239} % Extra information

\begin{document}

\maketitle

\chapter{Gravitational Field: a sphere, pt. 1}\label{quest: one}

\begin{exercise}{}
    Find the gravitational field inside and outside a sphere of radius $R$, which carries a uniform volume mass density $\rho$.
\end{exercise}

To find the gravitational field inside the sphere, we can use Gauss's law, which states that the gravitational flux through a closed surface is proportional to the enclosed mass. So, we can consider a spherical Gaussian surface of radius $r$ (where $r < R$) centered at the center of the sphere. The mass enclosed by this Gaussian surface is given by:
\begin{equation*}
    M_{\text{enclosed}} = \int\limits_{V} \rho(\vb{r}') \dd[3]{\vb{r}'} = \rho \int_{0}^{r} \int_{0}^{\pi} \int_{0}^{2\pi} (r')^{2} \sin\theta' \dd{\phi'} \dd{\theta'} \dd{r'} = \frac{4}{3} \pi r^{3} \rho
\end{equation*}

The gravitational field at a distance $r$ from the center of the sphere can be found using Gauss's law too:
\begin{equation*}
    \oint\limits_{S} \vb{g}(\vb{r}) \cdot \hat{\vb{n}} \dd{A} = -4\pi G M_{\text{enclosed}}
\end{equation*}

Since the gravitational field is radial and has the same magnitude at every point on the Gaussian surface, we can simplify this to:
\begin{equation*}
    g(r) 4\pi r^{2} = -4\pi G \qty(\frac{4}{3} \pi r^{3} \rho)
\end{equation*}

This gives us the gravitational field inside the sphere:
\begin{answer}
    \vb{g}(\vb{r}) = -\frac{4}{3} \pi G \rho r \hat{\vb{r}}, \quad \text{for } r < R
\end{answer}

To find the gravitational field outside the sphere, we can use the same argument, but now the Gaussian surface is a sphere of radius $r$ (where $r > R$) that encloses the entire mass of the sphere. Therefore, the enclosed mass is:
\begin{equation*}
    M_{\text{enclosed}} = \frac{4}{3} \pi R^{3} \rho
\end{equation*}

Using the same application of Gauss's law to find the gravitational field, we have:
\begin{equation*}
    g(r) 4\pi r^{2} = -4\pi G M_{\text{enclosed}} = -4\pi G \qty(\frac{4}{3} \pi R^{3} \rho)
\end{equation*}

This gives us the gravitational field outside the sphere:
\begin{answer}
    \vb{g}(\vb{r}) = -\frac{4}{3} \pi G \rho \frac{R^{3}}{r^{2}} \hat{\vb{r}}, \quad \text{for } r > R
\end{answer}

\endex

\newpage

\chapter{Gravitational Field: a spherical shell, pt. 1}\label{quest: two}

\begin{exercise}{}
    Find the gravitational field inside and outside a spherical shell of radius $R$, which carries a uniform surface mass density $\sigma$.
\end{exercise}

The gravitational field inside the spherical shell is zero, since there is no mass enclosed by any Gaussian surface inside the shell. Therefore, we have:
\begin{answer}
    \vb{g}(\vb{r}) = \vb{0}, \quad \text{for } r < R
\end{answer}

To find the gravitational field outside the spherical shell, we can use Gauss's law again. The mass enclosed by a Gaussian surface of radius $r$ (where $r > R$) is given by:
\begin{equation*}
    M_{\text{enclosed}} = \int\limits_{S} \sigma(\vb{r}') \dd[2]{\vb{r}'} = \sigma \int_{0}^{\pi} \int_{0}^{2\pi} R^{2} \sin\theta' \dd{\phi'} \dd{\theta'} = 4\pi R^{2} \sigma
\end{equation*}

We can use again that:
\begin{equation*}
    \oint\limits_{S} \vb{g}(\vb{r}') \cdot \dd[2]{\vb{r}'} = -4\pi G M_{\text{enclosed}}
\end{equation*}

Since the gravitational field is radial and has the same magnitude at every point on the Gaussian surface, we can simplify this to:
\begin{equation*}
    g(r) 4\pi r^{2} = -4\pi G M_{\text{enclosed}} = -4\pi G \qty(4\pi R^{2} \sigma)
\end{equation*}

This gives us the gravitational field outside the spherical shell:
\begin{answer}
    \vb{g}(\vb{r}) = -\frac{4\pi G R^{2} \sigma}{r^{2}} \hat{\vb{r}}, \quad \text{for } r > R
\end{answer}    

\newpage

\chapter{Gravitational potential: a straight wire}\label{quest: three}

\begin{exercise}{}
    Find the gravitational potential a distance $s$ from an infinitely long straight wire that carries a uniform line charge $\lambda$. Compute the gradient of your potential and find the corresponding gravitational field.
\end{exercise}

To find the gravitational potential at a distance $s$ from an infinitely long straight wire, we can use:
\begin{equation*}
    U(\vb{s}) = G \int \frac{\dd{m}}{\norm{\vb{s}}}
\end{equation*}

Since the wire is infinitely long, we can consider a small segment of the wire of length $\dd{z}$, which has a mass $\dd{m} = \lambda \dd{z}$. The distance from this segment to the point where we want to calculate the potential is given by $r = \sqrt{s^{2} + z^{2}}$. Taking a segment with length $2L$ centered at $z=0$, we can write:
\begin{equation*}
    U(s) = G \int_{-L}^{L} \frac{\lambda }{\sqrt{s^{2} + z^{2}}}\dd{z} = G \lambda \int_{-L}^{L} \frac{1}{\sqrt{s^{2} + z^{2}}}\dd{z}
\end{equation*}

The indefinite integral has a known solution:
\begin{equation*}
    \int \frac{1}{\sqrt{s^{2} + z^{2}}}\dd{z} = \ln\qty(z + \sqrt{s^{2} + z^{2}}) + C
\end{equation*}

So, the potential can be calculated as:
\begin{equation*}
    U(s) = G \lambda \qty[ \ln\qty(L + \sqrt{s^{2} + L^{2}}) - \ln\qty(-L + \sqrt{s^{2} + L^{2}}) ] = G \lambda \ln\qty(\frac{L + \sqrt{s^{2} + L^{2}}}{-L + \sqrt{s^{2} + L^{2}}})
\end{equation*}

The argument of the logarithm can be simplified as:
\begin{equation*}
    \frac{L + \sqrt{s^{2} + L^{2}}}{-L + \sqrt{s^{2} + L^{2}}} = \frac{L + \sqrt{s^{2} + L^{2}}}{ - L + \sqrt{s^{2} + L^{2}}} \cdot \frac{L + \sqrt{s^{2} + L^{2}}}{L + \sqrt{s^{2} + L^{2}}} = \frac{(L + \sqrt{s^{2} + L^{2}})^{2}}{s^{2}}
\end{equation*}

For large $L$, we can approximate $\sqrt{s^{2} + L^{2}} \approx L$, which gives us:
\begin{equation*}
    U(s) = 2G\lambda \ln\qty(\dfrac{2L}{s}) = 2G\lambda \ln(2L) - 2G\lambda \ln(s)
\end{equation*}

The first term diverges as $L \to \infty$, but does not depend on $s$, so it has no physical effect to the gravitational field. This can be absorbed into a reference potential, i.e., the difference:
\begin{equation*}
    U(s) - U(s_{0}) = -2G\lambda \ln(s) + 2G\lambda \ln(s_{0}) = -2G\lambda \ln\qty(\frac{s}{s_{0}})
\end{equation*}

Choosing $s_{0}$ such that $U(s_{0}) = 0$, we can write the potential as:
\begin{answer}
    U(s) = -2G\lambda \ln\qty(\dfrac{s}{s_{0}})
\end{answer}

To find the gravitational field, we can compute the gradient of the potential:
\begin{equation*}
    \vb{g}(\vb{s}) = \nabla U(s) = -2G\lambda \nabla \qty[\ln\qty(\dfrac{s}{s_{0}})] \hat{\vb{s}} = -2G\lambda \dv{}{s}\qty[\ln\qty(\dfrac{s}{s_{0}})] \hat{\vb{s}} = -2G\lambda \frac{1}{s} \hat{\vb{s}}
\end{equation*}

Concluding that:
\begin{answer}
    \vb{g}(\vb{s}) = -\frac{2G\lambda}{s} \hat{\vb{s}}
\end{answer}

\endex

\newpage

\chapter{Gravitational potential: a hollow sphere}\label{quest: four}

\begin{exercise}{}
    The gravitational potential $U(\theta)$ is specified on the surface of a hollow sphere of radius $R$. Find the potential inside the sphere.
\end{exercise}

Knowing that $\nabla\cdot \vb{g} = -4\pi G \rho$, we can use that $\vb{g} = \nabla U$ to obtain the Poisson equation for the gravitational potential $\nabla^{2} U = -4\pi G \rho$. In the case of a hollow sphere, we have $\rho = 0$ everywhere, so the Poisson equation reduces to the Laplace equation $\nabla^{2} U = 0$. The general solution to the Laplace equation in spherical coordinates can be written as:
\begin{equation*}
    U(r, \theta) = \sum_{\ell=0}^{\infty} \qty[A_{\ell} r^{\ell} + \frac{B_{\ell}}{r^{\ell+1}}] P_{\ell}(\cos\theta)
\end{equation*}

Inside the sphere, the solution must be finite at $r=0$, which implies that $B_{\ell} = 0$ for all $\ell$. Therefore, the potential inside the sphere can be written as:
\begin{equation*}
    U_{\text{inside}}(r, \theta) = \sum_{\ell=0}^{\infty} A_{\ell} r^{\ell} P_{\ell}(\cos\theta)
\end{equation*}

On the surface of the sphere, we have $U(R, \theta) = U(\theta)$, which gives us:
\begin{equation*}
    U(\theta) = \sum_{\ell=0}^{\infty} A_{\ell} R^{\ell} P_{\ell}(\cos\theta)
\end{equation*}

Multiplying both sides by $P_{\ell'}(\cos\theta)$ and integrating over $\dd{\cos\theta}$ from $-1$ to $1$:
\begin{equation*}
    \int_{-1}^{1} U(\theta) P_{\ell'}(\cos\theta)\dd{\cos\theta} = \sum_{\ell=0}^{\infty} A_{\ell} R^{\ell} \int_{-1}^{1} P_{\ell}(\cos\theta) P_{\ell'}(\cos\theta)\dd{\cos\theta}
\end{equation*}

Using the orthogonality of the Legendre polynomials, we have for the integral on the right-hand side the   result $\frac{2}{2\ell' + 1} \delta_{\ell \ell'}$, therefore:
\begin{equation*}
    \int_{-1}^{1} U(\theta) P_{\ell'}(\cos\theta)\dd{\cos\theta} = A_{\ell'} R^{\ell'} \frac{2}{2\ell' + 1}
\end{equation*}
implying that:
\begin{equation*}
    A_{\ell} = \frac{2\ell + 1}{2 R^{\ell}} \int_{-1}^{1} U(\theta) P_{\ell}(\cos\theta)\dd{\cos\theta}
\end{equation*}

Thus, the potential inside the sphere can be expressed as:
\begin{answer}
    U_{\text{inside}}(r, \theta) = \sum_{\ell=0}^{\infty} \frac{2\ell + 1}{2 R^{\ell}} r^{\ell} P_{\ell}(\cos\theta) \int_{-1}^{1} U(\theta) P_{\ell}(\cos\theta)\dd{\cos\theta} 
\end{answer}

\newpage

\chapter{Gravitational potential: a sphere, pt. 2}\label{quest: five}

\begin{exercise}{}
    The potential at the surface of a sphere of radius $R$ is given by $U(\theta) = k\cos(3\theta)$, where $k$ is a constant. Find the potential inside and outside the sphere, as well as the surface mass density $\sigma(\theta)$ on the sphere (assume there is no mass inside or outside the sphere).
\end{exercise}

By the Question 04, we know that we can write $U(r,\theta)$ in terms of the Legendre polynomials, therefore, we can express $U(\theta) = k\cos(3\theta)$ as a linear combination of some Legendre polynomials. Before we do that, the $\cos(3\theta)$ can be expanded as:
\begin{equation*}
    \cos(3\theta) = -3\cos\theta + 4\cos^{3}\theta
\end{equation*}

This means that we need Legendre polynomials of 1st and 3rd order in $\cos\theta$. The first few Legendre polynomials are:
\begin{align*}
    P_{0}(\cos\theta) &= 1 \\
    P_{1}(\cos\theta) &= \cos\theta \\
    P_{2}(\cos\theta) &= \frac{1}{2}(3\cos^{2}\theta - 1) \\
    P_{3}(\cos\theta) &= \frac{1}{2}(5\cos^{3}\theta - 3\cos\theta)
\end{align*}

In these cases, we can note that only $P_{1}(\cos\theta)$ and $P_{3}(\cos\theta)$ are needed to express $\cos(3\theta)$, therefore, we can write:
\begin{equation*}
    \cos(3\theta) = AP_{1}(\cos\theta) + BP_{3}(\cos\theta)
\end{equation*}
where $A$ and $B$ are constants to be determined. Substituting the expressions for $P_{1}(\cos\theta)$ and $P_{3}(\cos\theta)$, we have:
\begin{align*}
    \cos(3\theta) &\eq A\cos\theta + B\qty[\frac{1}{2}(5\cos^{3}\theta - 3\cos\theta)] \\
    &\eq \qty(A - \frac{3}{2}B)\cos\theta + \frac{5}{2}B\cos^{3}\theta
\end{align*}

Comparing the coefficients of the two expressions for $\cos(3\theta)$, we obtain:
\begin{equation*}
    \dfrac{5}{2}B = 4 \implies B = \frac{8}{5} \qquad \text{and} \qquad A - \frac{3}{2}B = -3 \implies A = -\frac{12}{5}
\end{equation*}

Thus, we have:
\begin{equation*}
    U(\theta) = \dfrac{2k}{5} \qty[-6P_{1}(\cos\theta) + 4P_{3}(\cos\theta)]
\end{equation*}

Knowing that $U(\theta)$ depends only on $P_{1}(\cos\theta)$ and $P_{3}(\cos\theta)$, the potential inside the sphere will be expressed only for $\ell = 1$ and $\ell = 3$, because of the orthogonality of the Legendre polynomials. Therefore, we can write:
\begin{equation*}
    U_{\text{inside}}(r, \theta) = A_{1} r P_{1}(\cos\theta) + A_{3} r^{3} P_{3}(\cos\theta)
\end{equation*}
where $A_{\ell}$ is given by \eqref{eq:Al}. Calculating $A_{1}$:
\begin{align*}
    A_{1} &\eq \dfrac{3}{2 R} \int_{-1}^{1} U(\theta) P_{1}(\cos\theta)\dd{\cos\theta} \\
    &\eq \dfrac{3}{2 R} \int_{-1}^{1} \dfrac{2k}{5} \qty[-6P_{1}(\cos\theta) + 4P_{3}(\cos\theta)] P_{1}(\cos\theta)\dd{\cos\theta} \\
    &\eq -\dfrac{36k}{10 R} \int_{-1}^{1} P_{1}^{2}(\cos\theta)\dd{\cos\theta} \\ 
\end{align*}

Using that:
\begin{equation*}
    \int_{-1}^{1} P_{\ell}^{2}(\cos\theta)\dd{\cos\theta} = \frac{2}{2\ell + 1}
\end{equation*}
we have:
\begin{equation*}
    A_{1} = -\dfrac{36k}{10 R} \cdot \dfrac{2}{3} = -\dfrac{12k}{5 R}
\end{equation*}

For $A_{3}$, we have:
\begin{align*}
    A_{3} &\eq \dfrac{7}{2 R^{3}} \int_{-1}^{1} U(\theta) P_{3}(\cos\theta)\dd{\cos\theta} \\
    &\eq \dfrac{7}{2 R^{3}} \int_{-1}^{1} \dfrac{2k}{5} \qty[-6P_{1}(\cos\theta) + 4P_{3}(\cos\theta)] P_{3}(\cos\theta)\dd{\cos\theta} \\
    &\eq \dfrac{56k}{10 R^{3}} \int_{-1}^{1} P_{3}^{2}(\cos\theta)\dd{\cos\theta} \\
    &\eq \dfrac{56k}{10 R^{3}} \cdot \dfrac{2}{7} = \dfrac{8k}{5 R^{3}}
\end{align*}

Thus, the potential inside the sphere can be expressed as:
\begin{align*}
    U_{\text{inside}}(r, \theta) &\eq -\dfrac{12k}{5 R} r P_{1}(\cos\theta) + \dfrac{8k}{5 R^{3}} r^{3} P_{3}(\cos\theta) \\
    &\eq -\dfrac{12k}{5 R} r \cos\theta + \dfrac{8k}{5 R^{3}} r^{3} \cdot \dfrac{1}{2}(5\cos^{3}\theta - 3\cos\theta) \\
    &\eq \dfrac{4k}{5 R} r \cos\theta \qty[-3 + \dfrac{r^{2}}{2R^{2}} (5\cos^{2}\theta - 3)]
\end{align*}

Concluding that the potential inside the sphere is given by:
\begin{answer}
    U_{\text{inside}}(r, \theta) = \dfrac{4k}{5 R} r \cos\theta \qty[-3 + \dfrac{r^{2}}{2R^{2}} (5\cos^{2}\theta - 3)]
\end{answer}

Outside the sphere, we can use the general solution to the Laplace equation to determine the potential. In this case, the solution must vanish at infinity, which implies that $A_{\ell} = 0$ for all $\ell$. Therefore, the potential outside the sphere can be written as:
\begin{equation*}
    U_{\text{outside}}(r, \theta) = \sum_{\ell=0}^{\infty} \frac{B_{\ell}}{r^{\ell+1}} P_{\ell}(\cos\theta)
\end{equation*}

On the surface of the sphere, we have $U(R, \theta) = U(\theta)$, which gives us:
\begin{equation*}
    U(\theta) = \sum_{\ell=0}^{\infty} \frac{B_{\ell}}{R^{\ell+1}} P_{\ell}(\cos\theta)
\end{equation*}

Multiplying both sides by $P_{\ell'}(\cos\theta)$ and integrating over $\dd{\cos\theta}$ from $-1$ to $1$:
\begin{align*}
    \int_{-1}^{1} U(\theta) P_{\ell'}(\cos\theta)\dd{\cos\theta} &\eq \sum_{\ell=0}^{\infty} \frac{B_{\ell}}{R^{\ell+1}} \int_{-1}^{1} P_{\ell}(\cos\theta) P_{\ell'}(\cos\theta)\dd{\cos\theta} \\
    &\eq \dfrac{2B_{\ell'}}{(2\ell' + 1) R^{\ell'+1}}
\end{align*}
implying that:
\begin{equation*}
    B_{\ell} = \frac{2\ell + 1}{2} R^{\ell+1} \int_{-1}^{1} U(\theta) P_{\ell}(\cos\theta)\dd{\cos\theta}
\end{equation*}

By ortogonality of the Legendre polynomials, only $B_{1}$ and $B_{3}$ will be non-zero, therefore:
\begin{align*}
    B_{1} &\eq \dfrac{3}{2} R^{2} \int_{-1}^{1} U(\theta) P_{1}(\cos\theta)\dd{\cos\theta} \\
    &\eq \dfrac{3}{2} R^{2} \int_{-1}^{1} \dfrac{2k}{5} \qty[-6P_{1}(\cos\theta) + 4P_{3}(\cos\theta)] P_{1}(\cos\theta)\dd{\cos\theta} \\
    &\eq -\dfrac{36k}{10} R^{2} \int_{-1}^{1} P_{1}^{2}(\cos\theta)\dd{\cos\theta} \\
    &\eq -\dfrac{36k}{10} R^{2} \cdot \dfrac{2}{3} \\
    &\eq -\dfrac{12k}{5} R^{2}
\end{align*}
and

\begin{align*}
    B_{3} &\eq \dfrac{7}{2} R^{4} \int_{-1}^{1} U(\theta) P_{3}(\cos\theta)\dd{\cos\theta} \\
    &\eq \dfrac{7}{2} R^{4} \int_{-1}^{1} \dfrac{2k}{5} \qty[-6P_{1}(\cos\theta) + 4P_{3}(\cos\theta)] P_{3}(\cos\theta)\dd{\cos\theta} \\
    &\eq \dfrac{56k}{10} R^{4} \int_{-1}^{1} P_{3}^{2}(\cos\theta)\dd{\cos\theta} \\
    &\eq \dfrac{56k}{10} R^{4} \cdot \dfrac{2}{7} \\
    &\eq \dfrac{8k}{5} R^{4}
\end{align*}

Thus, the potential outside the sphere can be expressed as:
\begin{align*}
    U_{\text{outside}}(r, \theta) &\eq -\dfrac{12k}{5} \frac{R^{2}}{r^{2}} P_{1}(\cos\theta) + \dfrac{8k}{5} \frac{R^{4}}{r^{4}} P_{3}(\cos\theta) \\
    &\eq -\dfrac{12k}{5} \frac{R^{2}}{r^{2}} \cos\theta + \dfrac{8k}{5} \frac{R^{4}}{r^{4}} \cdot \dfrac{1}{2}(5\cos^{3}\theta - 3\cos\theta) \\
    &\eq \dfrac{4k}{5} \frac{R^{2}}{r^{2}} \cos\theta \qty[-3 + \dfrac{R^{2}}{2r^{2}} (5\cos^{2}\theta - 3)]
\end{align*}

Concluding that the potential outside the sphere is given by:
\begin{answer}
    U_{\text{outside}}(r, \theta) = \dfrac{4k}{5} \frac{R^{2}}{r^{2}} \cos\theta \qty[-3 + \dfrac{R^{2}}{2r^{2}} (5\cos^{2}\theta - 3)]
\end{answer}

Because there is no mass inside or outside the sphere, the surface mass density $\sigma(\theta)$ can be obtained from the discontinuity of the normal component of the gravitational field across the surface of the sphere:
\begin{equation*}
    g_{\text{outside}}^{(r)} - g_{\text{inside}}^{(r)} = -4\pi G \sigma(\theta) \implies \sigma(\theta) = -\frac{1}{4\pi G} \qty[\pdv{U_{\text{outside}}}{r} - \pdv{U_{\text{inside}}}{r}]_{r=R}
\end{equation*}

For the derivative of $U_{\text{inside}}$ with respect to $r$, we have:
\begin{align*}
    \pdv{U_{\text{inside}}}{r} &\eq -\dfrac{12k}{5R} P_{1}(\cos\theta) + \dfrac{24k}{5R^{3}} r^{2} P_{3}(\cos\theta)
\end{align*}
and for the derivative of $U_{\text{outside}}$ with respect to $r$, we have:
\begin{align*}
    \pdv{U_{\text{outside}}}{r} &\eq \dfrac{24k}{5} \frac{R^{2}}{r^{3}} P_{1}(\cos\theta) - \dfrac{32k}{5} \frac{R^{4}}{r^{5}} P_{3}(\cos\theta)
\end{align*}

Thus:
\begin{align*}
    \qty[\pdv{U_{\text{outside}}}{r} - \pdv{U_{\text{inside}}}{r}]_{r=R} &\eq \dfrac{12k}{5R} P_{1}(\cos\theta) - \dfrac{8k}{5R} P_{3}(\cos\theta) \\
    &\eq \dfrac{4k}{5R} [3P_{1}(\cos\theta) - 2P_{3}(\cos\theta)] \\
    &\eq \dfrac{4k}{5R} \qty[3\cos\theta - (5\cos^{3}\theta - 3\cos\theta)] \\
    &\eq \dfrac{4k}{5R} (6\cos\theta - 5\cos^{3}\theta)
\end{align*}

Concluding that the surface mass density on the sphere is given by:
\begin{answer}
    \sigma(\theta) = -\dfrac{k}{5\pi G R} (6\cos\theta - 5\cos^{3}\theta)
\end{answer}

\endex

\newpage

\chapter{Gravitational potential: a spherical shell, pt. 2}\label{quest: six}

\begin{exercise}{}
    Derive the exact gravitational potential for a spherical shell of radius $R$ which carries a surface mass density $\sigma(\theta) = k\cos(\theta)$, where $k$ is a constant. Calculate also the quadrupole moment of this mass distribution.
\end{exercise}

Knowing the form of $\sigma(\theta)$, it can be written as $\sigma(\theta) = kP_{1}(\cos\theta)$. The general solutions of Laplace's equation give us:
\begin{equation*}
    U_{\text{inside}}(r, \theta) = \sum_{\ell=0}^{\infty} A_{\ell} r^{\ell} P_{\ell}(\cos\theta) \qquad \text{and} \qquad U_{\text{outside}}(r, \theta) = \sum_{\ell=0}^{\infty} \frac{B_{\ell}}{r^{\ell+1}} P_{\ell}(\cos\theta)
\end{equation*}

By the discontinuity of the normal component of the gravitational field across the shell, we have:
\begin{align*}
    \sigma(\theta) &\eq \frac{1}{4\pi G} \qty[\pdv{U_{\text{outside}}}{r} - \pdv{U_{\text{inside}}}{r}]_{r=R} \\
    &\eq \frac{1}{4\pi G} \qty[\sum_{\ell=0}^{\infty} A_{\ell} \ell R^{\ell-1} P_{\ell}(\cos\theta) + \sum_{\ell=0}^{\infty} (\ell + 1) \frac{B_{\ell}}{R^{\ell+2}} P_{\ell}(\cos\theta)]
\end{align*}

If $\sigma(\theta)$ depends only on $P_{1}(\cos\theta)$, then all terms in the above expression must vanish except for the one with $\ell = 1$, therefore:
\begin{align*}
    \sigma(\theta) &\eq \frac{1}{4\pi G} \qty[A_{1} R^{0} P_{1}(\cos\theta) + \frac{2B_{1}}{R^{3}} P_{1}(\cos\theta)] \\
    &\eq \frac{1}{4\pi G} \qty[A_{1} + \frac{2B_{1}}{R^{3}}] P_{1}(\cos\theta)
\end{align*}

We know that $U_{\text{inside}}(R, \theta) = U_{\text{outside}}(R, \theta)$, which implies that $A_{1} R P_{1}(\cos\theta) = \frac{B_{1}}{R^{2}} P_{1}(\cos\theta)$, therefore $B_{1} = A_{1} R^{3}$. Substituting this result in the above expression for $\sigma(\theta)$, we have:
\begin{align*}
    \sigma(\theta) &\eq \frac{1}{4\pi G} \qty[A_{1} + 2 A_{1}] P_{1}(\cos\theta) \\
    &\eq \frac{3 A_{1}}{4\pi G} P_{1}(\cos\theta)
\end{align*}

Thus:
\begin{equation*}
    \dfrac{3 A_{1}}{4\pi G} = k \implies A_{1} = \dfrac{4\pi G}{3} k \implies B_{1} = \dfrac{4\pi G}{3} k R^{3}
\end{equation*}

Therefore the potential inside the sphere can be expressed as:
\begin{answer}
    U_{\text{inside}}(r, \theta) = 4\pi G\dfrac{k r}{3} \cos\theta
\end{answer}
and the potential outside the sphere can be expressed as:
\begin{answer}
    U_{\text{outside}}(r, \theta) = 4\pi G \dfrac{kR^{3}}{3r^{2}} \cos\theta
\end{answer}

We can see that inside the sphere the potential is proportional to $r$, which give us a uniform gravitational field, while outside the sphere the potential is proportional to $\cos\theta/r^{2}$, which give us a dipole term in the multipole expansion of $U(\vb{r})$. This implies that the quadrupole moment of this mass distribution is zero, which can be verified by direct calculation:
\begin{align*}
    Q_{ij} &\eq \int \sigma(\theta') \qty(3 r'_{i} r'_{j} - r'^{2} \delta_{ij}) \dd[2]{r'} \\ 
    &\eq k \int \cos\theta' \qty(3 r'_{i} r'_{j} - r'^{2} \delta_{ij}) \dd[2]{r'} \\
    &\eq k R^{2} \qty(3 r'_{i} r'_{j} - \delta_{ij}) \int_{0}^{\pi} \int_{0}^{2\pi} \cos\theta'  \sin\theta' \dd{\phi'} \dd{\theta'} \\
    &\eq 2\pi k R^{2} \qty(3 r'_{i} r'_{j} - \delta_{ij}) \int_{0}^{\pi} \cos\theta' \sin\theta' \dd{\theta'} \\
    &\eq 0 (\text{by ortogonality})
\end{align*}

Concluding that the quadrupole moment of this mass distribution is zero:
\begin{answer}
    Q_{ij} = 0
\end{answer}

\endex

% \makefinal

\end{document}